


\section{Setting}
\label{sec:setting}


We distinguish three vector spaces and two maps that connect them.

\emph{Host parameters}
are $\haty \in \HostSpace\sim\mathbb{R}^{\nhost}$ quantities the analyst
wishes to differentiate with respect to.  
These are the ``outermost''
variables in the workflow---cross-section dimensions, load
multipliers, random-field coefficients, or any other inputs that the
user controls.

\emph{Core parameters} 
 $\corey \in \CoreSpace\sim\mathbb{R}^{\ncore}$ are quantities that
appear explicitly in the finite element residual and to which the DDM
can be applied directly.  
Typical core parameters include Young's
modulus, Poisson's ratio, fiber coordinates and areas, nodal
coordinates, and prescribed force magnitudes.

The \emph{model state} $\ModelState \in \mathbb{R}^{\nstate}$ is represented by the discrete solution
vector---for example, the vector of nodal displacements obtained by
solving the governing equilibrium equations
$\residual(\ModelState;\corey)=\mathbf{0}$.

\paragraph{Model map.}
The \emph{model map}
\begin{equation}\label{eq:Mmap}
    \Mmap : \mathbb{R}^{\ncore} \to \mathbb{R}^{\nstate},
    \qquad \corey \mapsto \ModelState = \Mmap(\corey),
\end{equation}
is the implicit function defined by the finite element analysis: given
core parameters~$\corey$, assemble the residual
$\residual(\ModelState;\corey)$, solve $\residual = \mathbf{0}$, and return
the equilibrium state~$\ModelState$.  The Jacobian
$\partial\Mmap/\partial\corey \in \mathbb{R}^{\nstate\times\ncore}$
is furnished column-by-column by the DDM; its $i$-th column is the
sensitivity vector $\partial\ModelState/\partial y_i$.

\paragraph{Parameter map.}
The \emph{parameter map}
\begin{equation}\label{eq:Ymap}
    \Ymap : \mathbb{R}^{\nhost} \to \mathbb{R}^{\ncore},
    \qquad \haty \mapsto \corey = \Ymap(\haty),
\end{equation}
is a user-defined, differentiable function that converts host
parameters into core parameters.  Because $\Ymap$ is implemented
entirely in the host language, its Jacobian
$\partial\Ymap/\partial\haty \in \mathbb{R}^{\ncore\times\nhost}$ is
available through automatic differentiation at no additional
implementation effort.


\paragraph{Composed map.}
The \emph{composed map} is the composition of the two:
\begin{equation}\label{eq:Tmap}
    \Tmap = \Mmap \circ \Ymap :
    \mathbb{R}^{\nhost} \to \mathbb{R}^{\nstate},
    \qquad \haty \mapsto \Mmap\!\bigl(\Ymap(\haty)\bigr).
\end{equation}
This is the end-to-end function that maps host parameters to the model
state.  The purpose of the framework is to make $\Tmap$ and its
derivatives available to gradient-based algorithms without requiring
the user to derive or implement the chain rule manually.

% Requires: \usepackage{tikz-cd}

\begin{equation}\label{eq:diagram}
\begin{tikzcd}[column sep=large, row sep=large]
    \mathbb{R}^{\nhost}
        \arrow[r, "\Ymap"]
        \arrow[dr, "\Tmap"']
    & \mathbb{R}^{\ncore}
        \arrow[d, "\Mmap"]
    \\
    & \mathbb{R}^{\nstate}
\end{tikzcd}
\end{equation}


\section{Differential Structure}
\label{sec:differential}


The chain rule for the composed map~\eqref{eq:Tmap} decomposes the
total derivative into a product of two Jacobians:
\begin{equation}\label{eq:chain-rule}
    \frac{\partial\Tmap}{\partial\haty}
    = \frac{\partial\Mmap}{\partial\corey}
      \;\frac{\partial\Ymap}{\partial\haty}
    \;\in\; \mathbb{R}^{\nstate \times \nhost}.
\end{equation}
In practice one never forms this matrix explicitly.  Instead, the
product is evaluated through either a \emph{Jacobian--vector product}
(JVP, forward mode) or a \emph{vector--Jacobian product} (VJP,
reverse mode), depending on the relative dimensions and the
application at hand.


\subsection{Forward mode: Jacobian--vector product}
% \label{sec:jvp}

Given a tangent vector $\dot{\haty} \in \mathbb{R}^{\nhost}$
representing a perturbation direction in the host parameter space, the
JVP of the composed map is
\begin{equation}\label{eq:jvp}
    \dot{\ModelState}
    = \mathrm{D}\Tmap(\haty)\;\dot{\haty}
    = \underbrace{%
        \mathrm{D}\Mmap\!\bigl(\Ymap(\haty)\bigr)
      }_{\text{kernel (DDM)}}
      \;\circ\;
      \underbrace{%
        \mathrm{D}\Ymap(\haty)\;\dot{\haty}
      }_{\text{host (AD forward)}}.
\end{equation}
The computation proceeds in two stages:
\begin{enumerate}
    \item \textbf{Host forward pass.}  Evaluate
          $\corey = \Ymap(\haty)$ and propagate the tangent through
          the parameter map to obtain
          $\dot{\corey} = \mathrm{D}\Ymap(\haty)\;\dot{\haty}$ using
          forward-mode AD.
    \item \textbf{Kernel forward pass.}  With the core parameters
          $\corey$ in hand, solve the primal problem
          $\residual(\ModelState;\corey)=\mathbf{0}$ and invoke the DDM to
          obtain the sensitivity vectors
          $\partial\ModelState/\partial y_i$ for each core parameter.
          Contract these with the tangent:
          \begin{equation}\label{eq:jvp-contract}
              \dot{\ModelState}
              = \sum_{i=1}^{\ncore}
                \frac{\partial\ModelState}{\partial y_i}\;\dot{y}_i.
          \end{equation}
\end{enumerate}

The cost of the JVP scales with~$\nhost$ (one tangent propagation per
host-parameter direction), making forward mode efficient when the
number of host parameters is small relative to the dimension of the
state.


\subsection{Reverse mode: vector--Jacobian product}
% \label{sec:vjp}


Given a cotangent vector $\bar{\ModelState} \in \mathbb{R}^{\nstate}$
(e.g.\ the gradient of an objective $J$ with respect to~$\ModelState$),
the VJP of the composed map is
\begin{equation}\label{eq:vjp}
    \bar{\haty}
    = \bigl[\mathrm{D}\Tmap(\haty)\bigr]^{\!\top} \bar{\ModelState}
    = \underbrace{%
        \bigl[\mathrm{D}\Ymap(\haty)\bigr]^{\!\top}
      }_{\text{host (AD reverse)}}
      \;\cdot\;
      \underbrace{%
        \bigl[\mathrm{D}\Mmap(\corey)\bigr]^{\!\top} \bar{\ModelState}
      }_{\text{kernel (DDM)}}.
\end{equation}
The computation now proceeds in the reverse order:
\begin{enumerate}
    \item \textbf{Kernel adjoint pass.}  Using the DDM sensitivities
          $\partial\ModelState/\partial y_i$ already available from the
          primal solve, compute the core-parameter cotangent
          \begin{equation}\label{eq:vjp-kernel}
              \bar{y}_i
              = \biggl(\frac{\partial\ModelState}{\partial y_i}\biggr)^{\!\top}
                \bar{\ModelState},
              \qquad i = 1,\ldots,\ncore.
          \end{equation}
    \item \textbf{Host reverse pass.}  Propagate
          $\bar{\corey} = (\bar{y}_1,\ldots,\bar{y}_{\ncore})^{\!\top}$
          back through the parameter map~$\Ymap$ using reverse-mode
          AD to obtain
          $\bar{\haty} = [\mathrm{D}\Ymap(\haty)]^{\!\top}\,\bar{\corey}$.
\end{enumerate}

The cost of the VJP is independent of~$\nhost$: a single reverse pass
through~$\Ymap$ delivers the gradient with respect to all host
parameters.  This makes reverse mode the natural choice whenever the
number of outputs (or objectives) is small relative to the number of
inputs.

\begin{remark}
    Because the DDM furnishes the full sensitivity matrix
    $\partial\ModelState/\partial\corey$ column-by-column, both the JVP
    and VJP of the model map~$\Mmap$ are available at the same
    computational cost.  The distinction between forward and reverse
    mode is therefore relevant only to the host map~$\Ymap$.
\end{remark}



% \section{Composition Protocol}
% \label{sec:protocol}


% The framework rests on a separation of responsibilities between
% three layers, each of which contributes one ingredient.

% \paragraph{Kernel primitive.}
% The finite element library exposes the model map~$\Mmap$ as a
% \emph{differentiable primitive}---an atomic operation whose primal
% evaluation, JVP rule, and VJP rule are each implemented internally
% using the DDM.  From the perspective of the AD framework, $\Mmap$ is
% a black box with known first-order differential structure.

% \paragraph{Host parameter map.}
% The user writes the parameter map~$\Ymap$ in the host language as
% ordinary, unrestricted code.  
% No special annotation or manual
% derivative is required; the AD framework differentiates~$\Ymap$
% automatically.

% \paragraph{Automatic composition.}
% The AD framework composes the two layers by the chain
% rule~\eqref{eq:chain-rule}.  From the user's perspective, the result
% is a single differentiable function
% $\Tmap : \haty \mapsto \ModelState$ that can be passed directly to any
% gradient-based algorithm---optimisation, uncertainty quantification,
% inverse problems---without further manual derivation.

% The key consequence of this decomposition is that extending the
% sensitivity analysis to new parameterisations requires changes only in
% the host code; the kernel primitive remains fixed.  Conversely,
% improvements to the finite element solver or the DDM implementation
% propagate automatically to every host-level workflow that uses the
% primitive.
